\begin{ouverture}[geometrie]{12}{Construire et comparer des triangles}{Partie 4 · Géométrie et espace}
{Cas d’égalité · Codages · Protocoles · Preuves}
\objectif{Construire un triangle avec des données suffisantes.}
\objectif{Reconnaître les trois cas d’égalité des triangles.}
\objectif{Rédiger une preuve en associant les sommets.}
\prerequis{Angles du triangle, inégalité triangulaire et médiatrice.}
\end{ouverture}
\pagestyle{otMain}\essentiel\fontsize{11}{15}\selectfont
\phantomsection\label{cours-12-1}
\essSection{1}{Les cas d’égalité}
\begin{definition}Deux triangles sont égaux s’ils peuvent se superposer : leurs côtés correspondants et leurs angles correspondants sont égaux. L’ordre des sommets doit respecter cette correspondance.\end{definition}
\begin{propriete}Trois critères suffisent : trois côtés respectivement égaux (CCC) ; deux côtés et l’angle compris respectivement égaux (CAC) ; un côté et les deux angles à ses extrémités respectivement égaux (ACA). Une figure retournée peut rester superposable.\end{propriete}
\label{fig-cours-12-1-cours-triangles-egaux}\figureGeo{cours-triangles-egaux}

\begin{attention}Deux longueurs seules ne fixent pas un triangle. Trois angles fixent une forme, mais pas une taille. Pour CAC, il faut bien l’angle situé entre les deux côtés connus ; un autre angle ne garantit pas le critère.\end{attention}
\clearpage
\phantomsection\label{cours-12-2}
\essSection{2}{Construire et démontrer}
\begin{savoirfaire}[Un protocole CCC]Pour $AB=6$, $AC=5$, $BC=4$ cm : tracer $[AB]$, puis le cercle de centre A et de rayon 5, puis celui de centre B et de rayon 4. Une intersection est C ; relier. L’autre intersection donne une figure symétrique, donc un triangle égal.\end{savoirfaire}
\begin{savoirfaire}[Vérifier l’existence]La plus grande longueur doit être strictement inférieure à la somme des deux autres pour un triangle non aplati. Pour une construction avec angles, leur somme doit être compatible avec les $180^\circ$ d’un triangle.\end{savoirfaire}
\begin{exemple}[Rédiger une preuve]On écrit les trois égalités de données, on nomme le cas utilisé et l’on conclut que les triangles sont égaux dans l’ordre choisi. On peut alors transférer une longueur ou un angle correspondant. Le dessin aide à chercher mais ne remplace pas ces raisons.\end{exemple}
\begin{savoirfaire}[Utiliser un outil]À la règle, au compas et au rapporteur, laisser les traits utiles. Sur un logiciel de géométrie, construire avec les contraintes (cercle, angle, parallèle), puis déplacer un point pour contrôler la stabilité. Une mesure visuelle seule ne démontre pas une propriété.\end{savoirfaire}

\clearpage\rubrique{Découverte}
\begin{blocLivre}{ex-12-1}\begin{exercice}[1]{Vérifier un triangle}
Peut-on construire un triangle de côtés 4, 5 et 6 cm ? Proposer la première étape.
\label{fig-ex-12-1}\figureGeo{tri-01}
\end{exercice}\end{blocLivre}
\begin{blocLivre}{ex-12-2}\begin{exercice}[1]{Un triangle impossible}
Peut-on construire un triangle non aplati de côtés 2, 3 et 6 cm ?
\label{fig-ex-12-2}\figureGeo{tri-02}
\end{exercice}\end{blocLivre}
\begin{blocLivre}{ex-12-3}\begin{exercice}[1]{Deux côtés et un angle}
Construire $AB=5$ cm, $AC=4$ cm et $\widehat{BAC}=60^\circ$. Nommer le type de données.
\label{fig-ex-12-3}\figureGeo{tri-03}
\end{exercice}\end{blocLivre}
\begin{blocLivre}{ex-12-4}\begin{exercice}[1]{Un côté et deux angles}
Construire $AB=6$ cm, $\widehat{BAC}=50^\circ$ et $\widehat{ABC}=60^\circ$. Calculer le troisième angle.
\label{fig-ex-12-4}\figureGeo{tri-04}
\end{exercice}\end{blocLivre}
\begin{blocLivre}{ex-12-5}\begin{exercice}[1]{Trois côtés correspondants}
On donne $AB=DE=4$, $AC=DF=5$, $BC=EF=6$ cm. Les triangles ABC et DEF sont-ils égaux ? Écrire la correspondance.
\label{fig-ex-12-5}\figureGeo{tri-05}
\end{exercice}\end{blocLivre}

\rubrique{Entraînement}
\begin{blocLivre}{ex-12-6}\begin{exercice}[2]{Un angle compris}
On donne $AB=DE=5$ cm, $AC=DF=3$ cm et $\widehat{BAC}=\widehat{EDF}=40^\circ$. Justifier l’égalité des triangles et en déduire une égalité de côtés.
\label{fig-ex-12-6}\figureGeo{tri-06}
\end{exercice}\end{blocLivre}
\begin{blocLivre}{ex-12-7}\begin{exercice}[2]{Deux angles aux extrémités}
$AB=DE=6$ cm, $\widehat A=\widehat D=50^\circ$ et $\widehat B=\widehat E=60^\circ$. Montrer que $\widehat C=\widehat F$ et que les triangles sont égaux.
\label{fig-ex-12-7}\figureGeo{tri-07}
\end{exercice}\end{blocLivre}
\begin{blocLivre}{ex-12-8}\begin{exercice}[2]{Des données insuffisantes}
Deux triangles ont chacun deux côtés de longueurs 5 et 4 cm. Peut-on conclure qu’ils sont égaux ?
\label{fig-ex-12-8}\figureGeo{tri-08}
\end{exercice}\end{blocLivre}
\begin{blocLivre}{ex-12-9}\begin{exercice}[2]{Un triangle isocèle}
ABC est isocèle en A, avec $AB=AC$ et $\widehat{BAC}=74^\circ$. Déterminer les deux autres angles.
\label{fig-ex-12-9}\figureGeo{tri-09}
\end{exercice}\end{blocLivre}
\begin{blocLivre}{ex-12-10}\begin{exercice}[2]{Un parallélogramme coupé}
ABCD est un parallélogramme et l’on trace la diagonale AC. Montrer que les triangles ABC et CDA sont égaux.
\label{fig-ex-12-10}\figureGeo{tri-10}
\end{exercice}\end{blocLivre}

\rubrique{Problèmes}
\begin{blocLivre}{ex-12-11}\begin{exercice}[3]{Un cerf-volant}
$AB=DB=5$ cm et $AC=DC=4$ cm. Les points A et D sont de part et d’autre de (BC). Montrer que $\widehat{BAC}=\widehat{BDC}$.
\label{fig-ex-12-11}\figureGeo{tri-11}
\end{exercice}\end{blocLivre}
\begin{blocLivre}{ex-12-12}\begin{exercice}[3]{Une hauteur et une aire}
Dans le triangle ABC, $AB=8$ cm et la hauteur CH vaut 3 cm. Calculer l’aire. Expliquer pourquoi AC ne peut pas remplacer CH dans la formule sans information supplémentaire.
\label{fig-ex-12-12}\figureGeo{tri-12}
\end{exercice}\end{blocLivre}
\begin{blocLivre}{ex-12-13}\begin{exercice}[3]{Un périmètre transféré}
Les triangles ABC et DEF sont égaux, avec A correspondant à D, B à E et C à F. On sait $AB=7$, $EF=5$ et $DF=4$ cm. Calculer le périmètre de ABC.
\label{fig-ex-12-13}\figureGeo{tri-13}
\end{exercice}\end{blocLivre}
\begin{blocLivre}{ex-12-14}\begin{exercice}[3]{Réparer un protocole}
On veut $AB=6$ cm, $AC=4$ cm et $\widehat{BAC}=35^\circ$. Un protocole place un angle de $35^\circ$ en B. Corriger et expliquer le problème.
\label{fig-ex-12-14}\figureGeo{tri-14}
\end{exercice}\end{blocLivre}
\begin{blocLivre}{ex-12-15}\begin{exercice}[3]{Démontrer une perpendicularité}
ABC est isocèle en A. M est le milieu de [BC]. Démontrer que (AM) est perpendiculaire à (BC) en utilisant l’égalité de deux triangles.
\label{fig-ex-12-15}\figureGeo{tri-15}
\end{exercice}\end{blocLivre}

\clearpage
